\documentclass[12pt,a4paper]{article}
\usepackage{amsmath,amsfonts,amssymb,amsthm}
\usepackage[OT1, OT2]{fontenc}
\usepackage[english,serbian]{babel}
\usepackage{paralist}
\usepackage{algorithm}
\usepackage{titlesec}

\topmargin       -11.0mm
\oddsidemargin      0mm
\evensidemargin     0mm
\textheight     235.5mm
\textwidth      161.0mm
\columnsep        4.1mm
\parindent        1.0em
\headsep          6.3mm
\headheight        12pt
\lineskip           1pt
\normallineskip     1pt
\def\baselinestretch{1.0}

\newcommand{\lN}{\mathbb{N}}
\newcommand{\lR}{\mathbb{R}}

\newtheorem{thm}{Teorema}%[section]
\renewcommand{\thethm}{\arabic{section}.\arabic{thm}}

\newtheorem{cor}{Korolar}%[section]
\renewcommand{\thethm}{\arabic{section}.\arabic{thm}}

\newtheorem{tvr}[thm]{Tvrdjenje}

\newtheorem{posl}[thm]{Posledica}

\newtheorem{defi}{Definicija}%[section]
\renewcommand{\thedefi}{\arabic{section}.\arabic{defi}}

\newtheorem*{nap}{Napomena}

\newtheorem{lem}{Lema}[subsection]

\newtheorem{pri}{Primer}[section]
\renewcommand{\thepri}{\arabic{section}.\arabic{pri}}

\newtheorem{zad}{Zadatak}

\newtheorem*{dokaz}{Dokaz}
%\renewcommand{\thedokaz}{}

\newcommand{\cyr}{\fontencoding{OT2}\selectfont\selectlanguage{serbian}}
\newcommand{\lat}{\fontencoding{T1}\selectfont\selectlanguage{english}}

\newtheorem{assumption}{Assumption}[section]
\renewcommand{\theassumption}{\arabic{section}.\arabic{assumption}}

\newtheorem{theorem}{Theorem}[section]
\renewcommand{\thetheorem}{\arabic{section}.\arabic{theorem}}

\newtheorem{lemma}{Lemma}[section]
\renewcommand{\thelemma}{\arabic{section}.\arabic{lemma}}

\newtheorem{corollary}{Corollary}[section]
\renewcommand{\thecorollary}{\arabic{section}.\arabic{corollary}}

\newtheorem{proposition}{Proposition}[section]
\renewcommand{\theproposition}{\arabic{section}.\arabic{proposition}}

\newcounter{appendixcounter}

\newtheorem{theoremappendix}{Theorem}
\renewcommand{\thetheoremappendix}{A\arabic{appendixcounter}.\arabic{theoremappendix}}

%\newtheorem{lemmaappendix}{Lemma}
%\renewcommand{\thelemmaappendix}{A\arabic{appendixcounter}.\arabic{lemmaappendix}}

\newtheorem{lemmaappendix}{Lemma}
\renewcommand{\thelemmaappendix}{A.\arabic{lemmaappendix}}

\newtheorem{assumptionappendix}{Assumption}[appendixcounter]
\renewcommand{\theassumptionappendix}{A\arabic{appendixcounter}.\arabic{assumptionappendix}}


\newtheorem*{remark}{Remark}%[section]
%\renewcommand{\theremark}{\arabic{section}.\arabic{remark}} 


\newtheorem{step}{Step}%[section]
\renewcommand{\thestep}{\arabic{step}} 

\titleformat{\section}[hang]{\bf\Large}{\thesection.}{0.4em}{}
\titleformat{\subsection}[hang]{\bf\large}{\thesubsection.}{0.4em}{}

\renewcommand{\baselinestretch}{1.01}

\begin{document}

\title{Tejlorova formula} 

\author{Milica V. Tadiћ
\thanks{Matamatichka Gimnazija, Kraljice Natalije 37, Beograd }}

\date{}

\maketitle

\section{Uvod} \label{section0}

U ovom seminarskom radu, razmatramo Tejlorovu formulu i njene osnovne osobine. 
Tejlorovu formulu mozhmo shvatiti/opisati kao 
lokalnu polinomijalnu aproksimaciju funkcija koje imaju izvode viseg reda. 
Formula je dobila ime po engleskom matematicharu 
Bruku Tejloru koji ju je formulisao u radu objavljenom 1715 godine. 

 
 
% \section{Tejlorova formula}
%Prirodno je postaviti pitanje mozhe li se funkcija $f$, koja %ima vishe izvoda u tachki
%a, ili njenoj okolini, aproksimirati polinomom visheg stepena %od jedan. Cilj
%te aproksimacije je da greshka bude shto je moguћe manja. %Odgovor na ovo pitanje daje
%Tejlorova formula.

\section{Tejlorov polinom i Tejlorov rezidual}

Prirodno se postavlja pitanje 
da li polinomom visheg stepena mozhe precizno 
aproksimirati 
funkciju koja ima vishe izvode u nekoj tachki i njenoj
okolini. 
Odgovor na ovo pitanje daje
Tejlorova formula. 
Tejlorova formula je bazirana 
na pojmu Tejlorovog polinoma i Tejlorovog reziduala. 


\begin{defi}[Tejlorov polinom i Tejlorov rezidual]
Neka je $f:\mathbb{R}\rightarrow\mathbb{R}$ funkcija 
koja je $n$ puta diferencijabilna u tachki $a\in\mathbb{R}$. 
Tejlorov polinom reda $n$ funkcije $f$ u tacki $a$ je 
polinom sledeћeg oblika 
\begin{align*}
    T_{n}(x,a)
    =
    \sum_{k=0}^{n}\frac{f^{(k)}}{k!}(x-a)^{k}, 
    \;\;\;\;\; x\in\mathbb{R}. 
\end{align*}
Odgovarajuћi Tejlorov rezidual je sledeћa funkcija
\begin{align*}
    R_{n}(x,a)
    =
    f(x)-T_{n}(x,a), 
    \;\;\;\;\; x\in\mathbb{R}. 
\end{align*}
\end{defi}

Korishћenjem Tejlorovog polinoma $T_{n}(x,a)$ i 
Tajlorovog reziduala $R_{n}(x,a)$, 
funkcija $f(x)$ se moze izraziti kao sledћi zbir 
\begin{align}\label{1.1}
    f(x)=T_{n}(x,a)+R_{n}(x,a). 
\end{align}
U matematichkoj analizi, jednachina (\ref{1.1}) je poznata 
kao Tejlorova formula. 

Ako je $x\approx a$, onda Tajlorov polinom $T_{n}(x,a)$ 
precizno aproksimira $f(x)$, tj. $T_{n}(x,a)\approx f(x)$. 
U tom slucaju, Tajlorov rezidual $R_{n}(x,a)$ predstavlja 
greshku aproksimacije. 
Da bih se procenila tachnost aproksimacije, 
potrebno je odrediti kako se Tajlorov rezidual 
ponasha kada se $x$ priblizhava $a$. 
Ovaj problem se 
razmatra u narednim odeljcima seminarksg rada. 

\begin{pri}
Neka je $f(x)=\exp(x)$ i $a=0$. 
Lako se onda proverava da je $f^{(n)}(a)=1$. 
Kao rezultat, dobijamo da je Tejlorov polinom %reda $n$
funkcije 
$f$ u tachki $a$ sledeћeg oblika
\begin{align*}
    T_{n}(x,a)
    =
    \sum_{k=0}^{n}\frac{x^{k}}{k!}. 
\end{align*}
\end{pri}

\begin{pri}
Neka je $f(x)=\cos(x)$ i $a=0$. 
Lako se onda proverava da je $f^{(2n)}(a)=(-1)^{n}$
i $f^{(2n+1)}(a)=0$. 
Za posledicu imamo da je Tejlorov 
polinom %redova $2n$ i $2n+1$ 
funkcije $f$ u tachki $a$ sledeћeg oblika
\begin{align*}
    T_{2n}(x,a)
    =
    T_{2n+1}(x,a)
    =
    \sum_{k=0}^{n}\frac{(-1)^{k}x^{2k}}{(2k)!}. 
\end{align*}
\end{pri}

\begin{pri}
Neka je $f(x)=\ln(x)$ i $a=1$. 
Lako se onda proverava da je $f^{(n)}(a)=(-1)^{n-1}(n-1)!$. 
Kao rezultat, dobijamo da je Tejlorov polinom %reda $n$
funkcije 
$f$ u tachki $a$ sledeћeg oblika
\begin{align*}
    T_{n}(x,a)
    =
    \sum_{k=0}^{n}\frac{(-1)^{k-1}(x-1)^{k}}{k}. 
\end{align*}
\end{pri}


\section{Rolova i Koshijeva teorema}

 U analizi Tejlorovog reziduala i njegovih osobina, 
 koriste se Rolova i Koshijeva teorema. 
 U ovom odeljku, formulisaћemo i dokazati ove 
 teoreme. 
 
 \begin{thm}[Rolova Teorema]\label{thm1}
 Neka su $a$ i $b$ realni brojevi za koje vazhi $a<b$. 
 Neka je $f:[a,b]\rightarrow\mathbb{R}$ funkcija 
 koja ispunjava sledeћe uslove: 
 \\
 (a) $f(a)=f(b)$. 
 \\
 (b) $f$ je neprekidna na intervalu $[a,b]$.
 \\
 (v) $f$ je diferencijabilna na intervalu 
 $(a,b)$. 
 \\
 Pod ovim uslovima, 
 postoji realni broj $c\in(a,b)$ tako da je $f'(c)=0$. 
 \end{thm}
 
 \begin{dokaz}
 Na osnovu Vajeshtrasove teoreme o ogranichenosti
 \cite[Teorema 14, Glava 5]{kadelburg1}, 
 postoje realni brojevi $\alpha$ i $\beta$ sadrzhani u 
 intervalu $[a,b]$ tako da je 
 \begin{align}\label{t1.1}
     f(\alpha)=\min_{x\in[a,b]}f(x), 
     \;\;\;\;\; 
     f(\beta)=\max_{x\in[a,b]}f(x). 
 \end{align}
 Kao neposrednu posledicu, dobijamo nejednakost
 \begin{align}\label{t1.3}
     f(\alpha)\leq f(x)\leq f(\beta)
 \end{align}
 za svako $x\in[a,b]$. 
 Iz (\ref{t1.3}) sledi $f(\alpha)\leq f(a)\leq f(\beta)$. 
 To za posledicu ima da 
 $a$, $\alpha$ i $\beta$ uvek ispunjavaju jedan od 
 sluchajeva $f(a)=f(\alpha)=f(\beta)$, $f(a)>f(\alpha)$ i 
 $f(a)<f(\beta)$. 
 U nastavku dokaza, ove sluchajeve odvojeno razmatramo. 
 \\
 {\bf\em Slucaj 1: $f(a)=f(\alpha)=f(\beta)$. }
 Koristeћi (\ref{t1.3}), zakljuchujemo da je 
 $f(x)=f(a)$ za svako $x\in[a,b]$. 
 Onda uslov (v) daje $f'(x)=0$ za svako $x\in(a,b)$. 
 To za posledicu ima da je $c\in(a,b)$ i $f'(c)=0$ za c=(a+b)/2. 
 \\
 {\bf\em Slucaj 2: $f(a)>f(\alpha)$. }
 Koristeћi uslov (a), zakljuchujemo da je $\alpha\in(a,b)$. 
 Istovremeno, (\ref{t1.3}) daje 
 \begin{align}\label{t1.5}
     \frac{f(\alpha+\delta)-f(\alpha)}{\delta}
     \leq 0, 
     \;\;\;\;\; 
     \frac{f(\alpha+\varepsilon)-f(\alpha)}{\varepsilon}
     \geq 0
 \end{align}
 za svako $\delta\in(a-\alpha,0)$ i $\varepsilon\in(0,b-\alpha)$. 
 Kombinujuћi uslov (v) i (\ref{t1.5}), zakljuchujemo 
 \begin{align*}
     f'(\alpha)
     =
     \lim_{h\rightarrow 0-}
     \frac{f(\alpha+h)-f(\alpha)}{h}
     \leq 0, 
     \;\;\;\;\; 
     f'(\alpha)
     =
     \lim_{h\rightarrow 0+}
     \frac{f(\alpha+h)-f(\alpha)}{h}
     \geq 0. 
 \end{align*}
 Kao neposrednu posledicu dobijamo $c\in(a,b)$ i $f'(c)=0$
 za $c=\alpha$. 
 \\
 {\bf\em Sluchaj 3: $f(a)<f(\beta)$. }
 Koristeћi uslov (a), zakljuchujemo da je $\beta\in(a,b)$. 
 Istovremeno, (\ref{t1.3}) daje 
 \begin{align}\label{t1.7}
     \frac{f(\beta+\delta)-f(\beta)}{\delta}
     \geq 0, 
     \;\;\;\;\; 
     \frac{f(\beta+\varepsilon)-f(\beta)}{\varepsilon}
     \leq 0
 \end{align}
 za svako $\delta\in(a-\beta,0)$ i $\varepsilon\in(0,b-\beta)$. 
 Kombinujuћi uslov (v) i (\ref{t1.5}), zakljuchujemo 
 \begin{align*}
     f'(\beta)
     =
     \lim_{h\rightarrow 0-}
     \frac{f(\beta+h)-f(\beta)}{h}
     \geq 0, 
     \;\;\;\;\; 
     f'(\beta)
     =
     \lim_{h\rightarrow 0+}
     \frac{f(\beta+h)-f(\beta)}{h}
     \leq 0. 
 \end{align*}
 Kao neposrednu posledicu dobijamo $c\in(a,b)$ i $f'(c)=0$
 za $c=\beta$. 
 \end{dokaz}
 
 
 \begin{thm}[Koshijeva Teorema]\label{thm2}
 Neka su $a$ i $b$ realni brojevi za koje vazhi $a<b$. 
 Neka su $f,g:[a,b]\rightarrow\mathbb{R}$ funkcije 
 koja ispunjavaju sledeћe uslove: 
 \\
 (a) $f$ i $g$ su neprekidne na intervalu $[a,b]$.
 \\
 (b) $f$ i $g$ su diferencijabilne na intervalu 
 $(a,b)$. 
 \\
 (v) $g'(x)\neq 0$ za svako $x\in(a,b)$. 
 \\
 Pod ovim uslovima, 
 postoji realni broj $c\in(a,b)$ tako da je 
 \begin{align}\label{t2.1*}
     \frac{f(b)-f(a)}{g(b)-g(a)}
     =
     \frac{f'(c)}{g'(c)}. 
 \end{align}
 \end{thm}
 
 \begin{nap}
 Koristeћi Rolovu teoremu i uslov (v), 
 zakljuchujemo da je $g(a)\neq g(b)$. 
 Naime, iz $g(a)=g(b)$ i Rolove teoreme direktno sledi 
 da je $g'(x)=0$ za neko $x\in(a,b)$ (a sto je nemoguћe zbog uslova (v)). 
 \end{nap}
 
 \begin{dokaz}
 Neka je $h:[a,b]\rightarrow\mathbb{R}$ funkcija sledeћeg oblika
 \begin{align}\label{t2.1}
     h(x)
     =
     f(x)
     -
     \frac{f(b)-f(a)}{g(b)-g(a)}g(x), 
     \;\;\;\;\; x\in[a,b]. 
 \end{align}
 Kao direktnu posledicu (\ref{t2.1}), dobijamo  
 \begin{align*}
     h(a)=h(b)
     =
     \frac{f(a)g(b)-f(b)g(a)}{g(b)-g(a)}. 
 \end{align*}
 Istovremeno, iz uslova (v) i (\ref{t2.1})
 proizilazi da je funkcija $h$ diferencijabilna na intervalu
 intervalu $(a,b)$. 
 Onda, primenjujuћi Rolovu teoremu na funkciju $h$, 
 zakljuchujemo da postoji $c\in(a,b)$ tako da je 
 $h'(c)=0$. 
 Iz uslova $h'(c)=0$ sledi 
 \begin{align*}
    f(c)
     -
     \frac{f(b)-f(a)}{g(b)-g(a)}g(c)
     =0, 
\end{align*}
a sto za direktnu posledicu ima (\ref{t2.1*}). 
 \end{dokaz}
 
 \subsection{Lagranzhov i Koshijev oblik Tejlorovog reziduala}
 
 U ovom odeljku, 
 pokazujemo da se Tejlorov rezidual mozhe predstaviti u 
 Lagranzhovom i Koshijevom obliku. 
 
 \begin{thm}\label{thm3}
 Neka su $a$ i $x$ realni brojevi za koje vazhi $a<x$. 
 Neka je $f:[a,x]\rightarrow\mathbb{R}$ funkcija 
 koja ispunjava sledeћe uslove: 
 \\
 (a) $f$ je $n$ puta diferencijabilna na intervalu 
 $[a,x]$. 
 \\
(b) $f^{(1)},\dots,f^{(n)}$ su neprekidni na intervalu 
 $[a,x]$. 
 \\
 (v) $f^{(n+1)}$ postoji na intervalu $(a,x)$.
 \\
 Pod ovim uslovima, postoje realni brojevi $u$ i $v$
 sadrzhani u intervalu $(a,x)$ tako da je 
 \begin{align}
    \label{t3.1*}
     R_{n}(x,a)
     =&
     \frac{f^{(n+1)}(u)}{(n+1)!}(x-a)^{n+1}, 
     \\
     \label{t3.3*}
     R_{n}(x,a)
     =&
     \frac{f^{(n+1)}(v)}{(n+1)!}(x-v)^{n}(x-a). 
 \end{align}
 \end{thm}
 
 \begin{nap}
 U matematichkoj analizi, jedancina (\ref{t3.1*}) je poznata 
 kao Lagranzhov oblik Tejlorovog reziduala, 
 dok Koshijev oblik istog reziduala dat u jednachini (\ref{t3.3*}). 
 U korolaru Teoreme \ref{thm3}, pokazacemo 
 da (\ref{t3.1*}) i (\ref{t3.3*}) vazhe i kada je $a>x$. 
 \end{nap}
 
 \begin{dokaz}
 Neka je $F,G,H:[a,x]\rightarrow\mathbb{R}$ funkcije sledћceg
 oblika 
 \begin{align}\label{t3.1}
     F(t)
     =
     \sum_{k=0}^{n}\frac{f^{(k)}(t)}{k!}(x-t)^{k}, 
     \;\;\;\;\; 
     G(t)=(x-t)^{n+1}, 
     \;\;\;\;\; 
     H(t)=t-a, 
     \;\;\;\;\; t\in[a,x]. 
 \end{align}
 Iz uslova (a) -- (b) i (\ref{t3.1}), zakljuchujemo da je 
 funkcija $F$ neprekidna na intervalu $[a,x]$ 
 i diferencijabilna na intervalu $(a,x)$. 
 Iz (\ref{t3.1}) takodje sledi da su funkcije $G$ i $H$ 
 neprekidne na intervalu $[a,x]$ 
 i diferencijabilne na intervalu $(a,x)$. 
 Onda, koristeћi Koshijevu teoremu, 
 zakljuchujemo da postoje realni brojevi $u$ i $v$ 
 sadrzhani u intervalu $(a,x)$ tako da je 
 \begin{align}\label{t3.3}
     \frac{F(x)-F(a)}{G(x)-G(a)}
     =
     \frac{F'(u)}{G'(u)}, 
     \;\;\;\;\; 
     \frac{F(x)-F(a)}{H(x)-H(a)}
     =
     \frac{F'(v)}{H'(v)}. 
 \end{align}
 
 Koristeћi prvu jednachinu u (\ref{t3.1}), dobijemo 
 \begin{align}\label{t3.5}
    F(x)-F(a)
    =
    f(x)-\sum_{k=0}^{n}\frac{f^{(k)}(a)}{k!}(x-a)^{k}
    =
    f(x)-T_{n}(x,a)
    =
    R_{n}(x,a).  
 \end{align}
 Iz iste jednachine, takodje dobijamo 
 \begin{align}\label{t3.7}
    F'(t)
    =&
    \sum_{k=0}^{n}
    \left(
    \frac{f^{(k+1)}(t)}{k!}(x-t)^{k}
    -
     \frac{f^{(k)}(t)}{k!}k(x-t)^{k-1}
    \right)
    \nonumber\\
    =&
    \sum_{k=0}^{n}
    \frac{f^{(k+1)}(t)}{k!}(x-t)^{k}
    -
    \sum_{k=1}^{n}
     \frac{f^{(k)}(t)}{(k-1)!}(x-t)^{k-1}
     \nonumber\\
     =&
     \frac{f^{(n+1)}(t)}{n!}(x-t)^{n}
 \end{align}
 za $t\in(a,x)$. 
 Istovremeno, iz druge jednachine 
 u (\ref{t3.1}) sledi  
 \begin{align}\label{t3.9}
    G(x)-G(a)=-(x-a)^{n+1}, 
    \;\;\;\;\; 
    G'(t)=-(n+1)(x-t)^{n},  
 \end{align}
 gde je $t\in(a,x)$. 
 Na slichan nachin, treћa jednachina u (\ref{t3.1}) daje 
 \begin{align}\label{t3.11}
     H(x)-H(a)=x-a, 
     \;\;\;\;\; 
     H'(t)=1, 
 \end{align}
 gde je $t\in(a,x)$. 
 
 Koristeћi (\ref{t3.5}), (\ref{t3.7}), (\ref{t3.9}) 
 i prvu jednachinu u (\ref{t3.3}), zakljuchujemo 
 \begin{align*}
     R_{n}(x,a)
     =
     F(x)-F(a)
     =
     \frac{F'(u)}{G'(u)}(G(x)-G(a))
     =
     \frac{f^{n+1}(u)}{(n+1)!}(x-a)^{n+1}. 
 \end{align*}
 Na slichan nachin, (\ref{t3.5}), (\ref{t3.7}), (\ref{t3.11}) 
 i druga jednachina u (\ref{t3.3}) daje  
 \begin{align*}
     R_{n}(x,a)
     =
     F(x)-F(a)
     =
     \frac{F'(v)}{H'(v)}(H(x)-H(a))
     =
     \frac{f^{n+1}(v)}{n!}(x-v)^{n}(x-a). 
 \end{align*}
 \end{dokaz}
 
 \begin{cor}\label{cor1}
 Neka su $a$ i $x$ realni brojevi za koje vazi $a>x$. 
 Neka je $f:[x,a]\rightarrow\mathbb{R}$ funkcija 
 koja ispunjava sledeћe uslove: 
 \\
 (a) $f$ je $n$ puta diferencijabilna na intervalu 
 $[x,a]$. 
 \\
(b) $f^{(1)},\dots,f^{(n)}$ su neprekidni na intervalu 
 $[x,a]$. 
 \\
 (v) $f^{(n+1)}$ postoji na intervalu $(x,a)$.
 \\
 Pod ovim uslovima, postoje realni brojevi $u$ i $v$
 sadrzhani u intervalu $(x,a)$ tako da
 (\ref{t3.1*}) i (\ref{t3.3*}) vaze. 
 \end{cor}
 
 \begin{dokaz}
 Neka je $\tilde{a}=-a$ i $\tilde{x}=-x$, 
 dok je $\tilde{f}:[\tilde{a},\tilde{x}]\rightarrow\mathbb{R}$ 
 funkcija oblika $\tilde{f}(t)=f(-t)$ za $t\in[a,b]$. 
 Onda, koristeћi uslove (a) -- (v), 
 zakljuchujemo da realni brojevi $\tilde{a}$, $\tilde{x}$
 i funkcija $\tilde{f}$ zadovoljavaju sve uslove 
 Teoreme \ref{t3.1}. 
 Ovo za posledicu ima da postoje realni brojevi 
 $\tilde{u}$ i $\tilde{v}$ sadrzhani u $(\tilde{a},\tilde{x})$
 tako da je 
 \begin{align*}
     \tilde{f}(\tilde{x})
     -
     \sum_{k=0}^{n}
     \frac{\tilde{f}^{(k)}(\tilde{a})}{k!}
     (\tilde{x}-\tilde{a})^{k}
     =&
     \frac{\tilde{f}^{n+1}(\tilde{u})}{(n+1)!}
     (\tilde{x}-\tilde{a})^{n+1}, 
     \\
     \tilde{f}(\tilde{x})
     -
     \sum_{k=0}^{n}
     \frac{\tilde{f}^{(k)}(\tilde{a})}{k!}
     (\tilde{x}-\tilde{a})^{k}
     =&
     \frac{\tilde{f}^{n+1}(\tilde{v})}{(n+1)!}
     (\tilde{v}-\tilde{a})^{n}(\tilde{x}-\tilde{a}). 
 \end{align*}
 
 Neka je $u=-\tilde{u}$ i $v=-\tilde{v}$. 
 Lako se onda zakljuchuje da su $u$ i $v$ sadrzhani u $(x,a)$
 i da ispunjavaju (\ref{t3.1}), (\ref{t3.3}). 
 \end{dokaz}
 
 
 \section{Tejlorova teorema}
 
 U ovom odeljku razmatramo asimptotsko ponasanje 
 Tejlorovog reziduala kada njegov prvi argument 
 tezhi drugom argumentu. 
 Odgovor na ovaj problem daje Tejlorova teorema. 
 
 \begin{thm}[Tejlorova Teorema]\label{thm4}
 Neka su $a$ i $b$ realni brojevi za koje vazhi $a<b$. 
 Neka je $f:[a,b)\rightarrow\mathbb{R}$ funkcija 
 koja ispunjava sledeћe uslove: 
 \\
 (a) $f$ je $n$ puta diferencijabilna na intervalu 
 $[a,b)$. 
 \\
(b) $f^{(1)},\dots,f^{(n)}$ su neprekidni na intervalu 
 $[a,b)$. 
 \\
 Pod ovim uslovima imamo da je 
 $R_{n}(x,a)=o\left((x-a)^{n}\right)$ kada $x\rightarrow a$, 
 tj., 
 \begin{align}
     \lim_{x\rightarrow a}
     \frac{R_{n}(x,a)}{(x-a)^{n}}=0. 
 \end{align}
 \end{thm}
 
 \begin{dokaz}
Na osnovu Teoreme \ref{thm3}, za svaki realni broj 
$x\in[a,b)$ postoji realni broj $u=u(x)\in(a,x)$ 
tako da je 
\begin{align}
    R_{n-1}(x,a)
    =
    \frac{f^{(n)}(u)}{(n)!}(x-a)^{n}. 
\end{align}
Za neposrednu posledicu imamo da 
$u=u(x)\rightarrow a$ kada $x\rightarrow a$. 
Takodje imamo 
\begin{align}
    \frac{R_{n}(x,a)}{(x-a)^{n}}
    =
    \frac{1}{(x-a)^{n}}
    \left(
    R_{n-1}(x,a)
    -
    \frac{f^{(n)}(a)}{n!}(x-a)^{n}
    \right)
    =
    \frac{f^{(n)}(u)-f^{(n)}(a)}{n!}. 
\end{align}
Kako je $f^{(n)}$ neprekidna na $[a,b)$, dobijamo 
\begin{align*}
    \lim_{x\rightarrow a}
    \frac{R_{n}(x,a)}{(x-a)^{n}}
    =
    \lim_{u\rightarrow a}
    \frac{f^{(n)}(u)-f^{(n)}(a)}{n!}
    =0. 
\end{align*}
\end{dokaz}


\section{Tejlorov red}

U matematichkoj analizi, 
znachajnu ulogu igraju funkcije koje su beskonachno 
mnogo puta 
diferencijabilne, tj. koje imaju sve izvode visheg reda. 
Takve funcije se nazivaju glatke funkcije. 
Jedna od osnovnih osobina glatkih funkcija je 
da imaju Tejlorov polinom bilo kog reda. 
U tom sluchaju, prirodno se nameћe pitanje 
da li Tejlorov polinom konvergira, i ako konvergira, 
pod kojim uslovima i kojoj granichnoj vrednosti. 
Odgovor na ovo pitanje bazirano je na pojmu 
Tejlorovog reda. 

\begin{defi}[Tejlorov red]
Neka je $f:\mathbb{R}\rightarrow\mathbb{R}$ funkcija 
koja je beskonachno mnogo puta diferencijabilna u 
tachki $a$. 
Tejlorov red funkcije $f$ u tachki $a$ je red 
sledeћeg oblika
\begin{align*}
    T_{\infty}(x,a)
    =
    \sum_{k=0}^{\infty}
    \frac{f^{(k)}(a)}{k!}(x-a)^{k}, 
    \;\;\;\;\; x\in\mathbb{R}. 
\end{align*}
\end{defi}

Neke od osnovnih osobina 
Tejlorovog reda sadrzhane su u sledeћoj teoremi. 

\begin{thm}\label{thm5}
Neka je $a$ realan broj, 
dok je $f:\mathbb{R}\rightarrow\mathbb{R}$ funkcija koja 
ispunjava sledeћe uslove: 
\\
(a) Postoji positivan realan broj $r$
tako da je $f$ beskonachno mnogo puta differencijabilna 
na intervalu $(a-r,a+r)$. 
\\
(b) Postoji pozitivan realan broj $C$ 
tako da je 
$|f^{(n)}(x)|\leq C^{n}n!$ za svako 
$x\in(a-r,a+r)$ i $n\geq 0$. 
\\
Pod ovim uslovima, 
Tejlorov red $T_{\infty}(x,a)$ je apsolutno kovergentan 
za svako $x\in(a-\rho,a+\rho)$, 
gde je $\rho=\min\{r,1/C\}$. 
Pod istim uslovima, imamo i da je 
\begin{align}
    f(x)=T_{\infty}(x,a)
    =
    \lim_{n\rightarrow\infty}T_{n}(x,a), 
    \;\;\;\;\; 
    |f(x)-T_{n}(x,a)|
    \leq 
    (C\rho)^{n+1}
\end{align}
za svako $x\in(x-\rho,x+\rho)$. 
\end{thm}

\begin{dokaz}
Neka je $x$ bilo koji realni broj u intervalu 
$(a-\rho,a+\rho)$. 
Kako je $\rho<r$ i $C\rho<1$, uslov (b) daje 
\begin{align*}
    \sum_{n=0}^{\infty}
    \frac{|f^{(n)}(a)|}{n!}|x-a|^{n}
    \leq 
    \sum_{n=0}^{\infty}
    (C\rho)^{n}
    =
    \frac{1}{1-Cr}<\infty. 
\end{align*}
To za posledicu ima da je 
Tejlorov red $T_{\infty}(x,a)$ apsolutno 
kovergentan i da limit 
$\lim_{n\rightarrow\infty}T_{n}(x,a)=T_{\infty}(x,a)$
vazhi. 
Istovremeno, na osnovu Teoreme \ref{thm3} i Korolara \ref{cor1}, 
za svako $n\geq 0$ postoji realna broj 
$u_{n}\in(a-r,a+r)$ tako da je 
\begin{align}\label{t5.1}
    R_{n}(x,a)
    =
    \frac{f^{(n+1)}(u_{n})}{(n+1)!}(x-a)^{n+1}. 
\end{align}
Koristeћi uslov (b) i (\ref{t5.1}), dobijamo 
\begin{align*}
    |f(x)-T_{n}(x,a)|
    =
    |R_{n}(x,a)|
    \leq 
    \frac{|f^{(n+1)}(u_{n})|}{(n+1)!}|x-a|^{n+1}
    \leq 
    (C\rho)^{n+1}. 
\end{align*}
Kako je $\lim_{n\rightarrow\infty}(C\rho)^{n}=0$, 
zakljuchujemo da je $f(x)=\lim_{n\rightarrow\infty}T_{n}(x,a)$. 
Kao neposrednu posledicu, dobijamo $f(x)=T_{\infty}(x,a)$. 
\end{dokaz}

\begin{pri}
Neka je $f(x)=\exp(x)$ i $a=0$. 
Kako je $f^{(n)}(a)=1$,  
Tejlorov red 
funkcije $f$ u tachki $a$ je sledeћeg oblika
\begin{align*}
    T_{\infty}(x,a)
    =
    \sum_{k=0}^{\infty}\frac{x^{k}}{k!}. 
\end{align*}
Istovremeno, imamo da je su izvodi $|f^{(n)}(x)|$ 
nezavisni od $n$. 
Onda, koristeћi Teoremu \ref{thm5}, 
zakljuchujemo da je $T_{\infty}(x,a)$ apsolutno konvergentan
za svako $x\in\mathbb{R}$
i da $f(x)=T_{\infty}(x,a)$ vazhi za isto $x$. 
\end{pri}

\begin{pri}
Neka je $f(x)=\cos(x)$ i $a=0$. 
Kako je $f^{(2n)}(a)=(-1)^{n}$
i $f^{(2n+1)}(a)=0$,  
Tejlorov red 
funkcije $f$ u tachki $a$ je sledeћeg oblika
\begin{align*}
    T_{\infty}(x,a)
     =
    \sum_{k=0}^{\infty}\frac{(-1)^{k}x^{2k}}{(2k)!}. 
\end{align*}
Istovremeno, imamo da je $|f^{(n)}(x)|\leq 1$ 
za svako $x\in\mathbb{R}$. 
Onda, koristeћi Teoremu \ref{thm5}, 
zakljuchujemo da je $T_{\infty}(x,a)$ apsolutno konvergentan
za svako $x\in\mathbb{R}$
i da $f(x)=T_{\infty}(x,a)$ vazhi za isto $x$. 
\end{pri}

\begin{pri}
Neka je $f(x)=\ln(x)$ i $a=1$. 
Kako je $f^{(n)}(a)=(-1)^{n-1}(n-1)!$, 
Tejlorov red
funkcije $f$ u tachki $a$ je sledeћeg oblika
\begin{align*}
    T_{\infty}(x,a)
    =
    \sum_{k=0}^{\infty}\frac{(-1)^{k-1}(x-1)^{k}}{k}. 
\end{align*}
Istovremeno, imamo da je $|f^{(n)}(x)|\leq (n-1)!/(1-r)^{n}$ 
za svako $x\in(1-r,1+r)$ i za svako $r\in(0,1)$. 
Onda, koristeћi Teoremu \ref{thm5}, 
zakljuchujemo da je $T_{\infty}(x,a)$ apsolutno konvergentan
za svako $x\in(0,2)$
i da $f(x)=T_{\infty}(x,a)$ vazhi za isto $x$. 
\end{pri}

\begin{thebibliography}{10}

\bibitem{kadelburg1}
Zoran Kadelburg, Vladimir Miћiћ, Srdjan Ognjanoviћ 
i Sonja Chukiћ,  
{\em Analiza sa Algebrom, Udџbenik sa zbirkom zadataka za 3. razred Matematichke gimnazije,}
Krug, Beograd, 2018. 
\cyr
\bibitem{kadelburg2}
Zoran Kadelburg, Vladimir Miћiћ и Srdjan Ognjanoviћ ,  
{\em Analiza sa Algebrom, Udџbenik sa zbirkom zadataka za 4. razred Matematichke gimnazije,}
Krug, Beograd, 2016. 
\end{thebibliography}

\end{document}
